% page 116 du fac-simile (image p-115 du scan)
% bloc 162.6 x 237.7 mm ; marge gauche 8.0 mm
\pgc{-12.700mm}{0.000mm}{
\l{\cel{7}lO6.}
\l{}
\l{\sou{diademo}. I.Bendo ornizita ye broduri, lapidi, quan la suvereni}
\l{\cel{3}+weris zone sur sua fronto. II. Bendo, kapo-ringa, ye qua}
\l{\cel{3}konsistas la krono di la suvereni. - DEFIRS.}
\l{}
\l{\sou{diafana}. Qua lasas pasar la lumo-radii, tale ke onu povas}
\l{\cel{3}dicernar la objekti tra lu. - EFIS.}
\l{}
\l{\sou{diafizo}. (\sou{anat.}) Korpo di la osti longa. - DEF.}
\l{}
\l{\sou{diafragmo}. I. (\sou{anat.})\cel{2}Muskulo parieta inter la pektoro e}
\l{\cel{3}la abdomino. - Parieto qua dividas ula organi. - II.(Op-}
\l{\cel{3}\sou{tiko}). Cirklo opaka, muntita en teleskopo o lorno, por}
\l{\cel{3}haltigar la radii qui ne kuniras a la foko. - Disko kun}
\l{\cel{3}truo ronda, quan la fotografisto muntas en la objektalo}
\l{\cel{3}samaskope. - EFIRS.}
\l{}
\l{\sou{diagnozar}.\cel{2}(\sou{trans.}) Determinar (morbo) per la simptomi qui}
\l{\cel{3}esas lua propraji. - DEFIRS.}
\l{}
\l{\sou{diagonala}. Qua iras de angulo ad angulo opozita. - DEFIRS.}
\l{}
\l{\sou{diagramo}. Expresuro grafika, tabelo lineala, cifrala, ek la}
\l{\cel{3}parti di ensemblo ed ek lia inter-relati. - DEFIRS.}
\l{}
\l{\sou{diakaustika}. (\sou{geom.}) Produktita da la inter-seki di la radii}
\l{\cel{3}refraktita.}
\l{}
\l{\sou{diakilono}. I. Plastro ek dekokturo de gladiolo-radiki, de oleo}
\l{\cel{3}e de litargiro, a qui onu adjuntas ulakaze gumo-rezini,}
\l{\cel{3}pecho blanka, terpentin-oleo, vaxo. - II. Sparadrapo ek ta}
\l{\cel{3}plastro, extensita sur telo peco. - DEF.}
\l{}
\l{\sou{diakono}. (en\cel{1}\sou{la eklezio katolik}a). Kleriko qua havas la duesma}
\l{\cel{3}ek la ordi superiora. - DEFIRS.}
\l{}
\l{\sou{diakoniso}. (\sou{religio)}.\cel{2}I. En la komenco-periodo di la eklezio}
\l{\cel{3}Kristana : filiino o vidvino qua recevis la impozo di la}
\l{\cel{3}manui, e su konsakris a la entrateno di la templi, di la}
\l{\cel{3}povri, e c. - II. En ula sekti di protestantismo : damo qua}
\l{\cel{3}su konsakris a verki piesala, karitatala. - DEFIRS.}
\l{}
\l{\sou{diakritika}.\cel{2}Qua esas signo adjuntita a litero (supersigno, sub-}
\l{\cel{3}signo, e c.) - kom exempli : ç, â, ĝ, ë,\cel{2}e c.) - DEFIRS.}
\l{}
\l{\sou{dialekto}. (\sou{linguistiko)}.\cel{2}Varietato regionala di linguo. -}
\l{\cel{3}DEFIRS.}
\l{}
\l{\sou{dialektiko}. Metodo per qua onu deduktas rezoni, uzata por}
\l{\cel{3}demonstrar o refutar. - DEFIRS.}
\l{}
\l{\sou{dializar}. (\sou{trans.})\cel{2}Separar du substanci dissolvura, per di-}
\l{\cel{3}fuzo tra parieto poroza, quan un de li povas trairar. -}
\l{\cel{3}DEF.}
}
